%% CAPA — 250-word conference abstract
%% Use for: AMIA High School Scholars, BMES Poster Expo, JSHS, VJAS, Regeneron STS
%%
%% Title (short form for conference headers):
%%   CAPA: Protein Language Model Embeddings for Competing-Risks Outcome
%%   Prediction in Pediatric Hematopoietic Stem Cell Transplantation
%%
%% Title (full form for paper / Regeneron STS):
%%   CAPA: Structure-Aware HLA Mismatch Representations with Protein Language
%%   Models for Competing-Risks Survival Prediction in Paediatric Haematopoietic
%%   Stem Cell Transplantation

\documentclass[11pt,letterpaper]{article}
\usepackage[margin=2.5cm]{geometry}
\usepackage[T1]{fontenc}
\usepackage[utf8]{inputenc}
\usepackage{lmodern}
\usepackage{amsmath}
\usepackage{hyperref}
\usepackage{microtype}

\title{\textbf{CAPA: Protein Language Model Embeddings for\\
Competing-Risks Outcome Prediction in Pediatric\\
Hematopoietic Stem Cell Transplantation}}
\author{Huanxuan (Shawn) Li\\
\small Thomas Jefferson High School for Science and Technology\\
\small \href{mailto:2027hli@tjhsst.edu}{2027hli@tjhsst.edu}}
\date{}

\begin{document}
\maketitle

%% ── Abstract (250 words) ──────────────────────────────────────────────────────

\begin{abstract}

\noindent\textbf{Background.}
Allogeneic hematopoietic stem cell transplantation (HSCT) is curative for
severe hematological malignancies, yet post-transplant outcomes are shaped by
immunological interactions that categorical HLA mismatch scoring cannot fully
capture.  Clinical practice treats all mismatches as equivalent regardless of
amino-acid-level structural divergence, and standard survival models analyze
graft-versus-host disease (GvHD), relapse, and transplant-related mortality
(TRM) as independent events---despite being mutually competing.

\medskip
\noindent\textbf{Methods.}
We present CAPA (Computational Architecture for Predicting Alloimmunity), a
deep learning framework for pediatric HSCT outcome modeling.  Each HLA allele
is encoded as a 1280-dimensional vector via the ESM-2 650M protein language
model.  A bidirectional cross-attention network extracts donor--recipient
interaction features, which are concatenated with clinical covariates and
passed to a DeepHit survival head that jointly estimates cumulative incidence
functions (CIFs) for GvHD, relapse, and TRM over a 730-day horizon.
We evaluate competing-risks baselines---Fine--Gray subdistribution hazard and
cause-specific Cox models---on the public UCI Bone Marrow Transplant dataset
($n=187$ pediatric patients; 70/15/15 train/val/test split).

\medskip
\noindent\textbf{Results.}
On the held-out test set ($n=29$), the Fine--Gray model reaches concordance
indices of 0.84 (95\,\% CI 0.69--1.00) for relapse and 0.66 (0.48--0.86) for
TRM; the cause-specific Cox model achieves 0.75 and 0.65, respectively.  A
flat-feature DeepHit MLP performs below classical baselines, consistent with
known training challenges on small cohorts.

\medskip
\noindent\textbf{Conclusions.}
CAPA demonstrates a protein language model framework for structure-aware HLA
mismatch representation in HSCT outcome modeling.  Full validation of ESM-2
embeddings requires registry-scale allele-level data, identified as the primary
direction for future work.  Source code and a multi-donor web interface are
publicly available at \url{https://github.com/sh4wn27/capa}.

\end{abstract}

\end{document}
